% CSE 711 — Compiler
% All Non-Numerical (Theory / Conceptual) Questions — Complete Model Answers (2016–2024)
% University of Chittagong, Graduate Exam
%
% Scope: definitions, comparisons and explanation questions across all 8 topics.
% Excludes construction/trace/computation problems (derivations, FIRST/FOLLOW,
% LL(1)/LR table construction, parse traces, annotated trees, DAGs, TAC generation,
% activation-record diagrams, instruction-cost calculation).
%
% Compile with: pdflatex Compiler_Theory_Solutions.tex (run twice for TOC)

\documentclass[12pt, a4paper]{article}

\usepackage[left=2.3cm, right=2.3cm, top=2.6cm, bottom=2.6cm, headheight=15pt]{geometry}
\usepackage{amsmath, amssymb}
\usepackage{booktabs}
\usepackage{array}
\usepackage{xcolor}
\usepackage{enumitem}
\usepackage{fancyhdr}
\usepackage{titlesec}
\usepackage{mdframed}
\usepackage{tabularx}
\usepackage{multirow}
\usepackage{hyperref}

%----- Colours -----
\definecolor{sectionblue}{RGB}{10,60,110}
\definecolor{boxbg}{RGB}{242,247,255}
\definecolor{answerbg}{RGB}{240,255,240}
\definecolor{notesbg}{RGB}{255,250,230}
\definecolor{keyblue}{RGB}{0,80,160}
\definecolor{accent}{RGB}{180,40,40}

%----- Box environments -----
\mdfdefinestyle{qframe}{linecolor=sectionblue, backgroundcolor=boxbg, roundcorner=4pt, linewidth=1.2pt, innertopmargin=6pt, innerbottommargin=6pt, innerleftmargin=8pt, innerrightmargin=8pt}
\mdfdefinestyle{notesframe}{linecolor=orange!70!black, backgroundcolor=notesbg, roundcorner=4pt, linewidth=1pt, innertopmargin=5pt, innerbottommargin=5pt, innerleftmargin=8pt, innerrightmargin=8pt}
\newenvironment{qbox}{\begin{mdframed}[style=qframe]}{\end{mdframed}}
\newenvironment{notesbox}{\begin{mdframed}[style=notesframe]}{\end{mdframed}}

%----- Page style -----
\pagestyle{fancy}
\fancyhf{}
\lhead{\small\color{sectionblue} CSE 711 — Compiler}
\rhead{\small\color{sectionblue} Theory Model Answers (2016–2024)}
\cfoot{\small\thepage}
\renewcommand{\headrulewidth}{0.5pt}

%----- Section formatting -----
\titleformat{\section}{\Large\bfseries\color{sectionblue}}{}{0em}{}[\vspace{-2pt}\color{sectionblue}\rule{\linewidth}{0.8pt}]
\titleformat{\subsection}{\large\bfseries\color{sectionblue!85!black}}{}{0em}{}

%----- Custom lists -----
\setlist[itemize]{leftmargin=1.4em, itemsep=1.5pt, topsep=2pt}
\setlist[enumerate]{leftmargin=1.6em, itemsep=2pt, topsep=2pt}

%----- Macros -----
\newcommand{\yr}[1]{{\small\color{accent}\textbf{[#1]}}}
\newcommand{\term}[1]{\textbf{\color{keyblue}#1}}
\renewcommand{\arraystretch}{1.25}
\newcolumntype{Y}{>{\raggedright\arraybackslash}X}

\hypersetup{colorlinks=true, linkcolor=sectionblue, urlcolor=sectionblue}

\begin{document}

%==================== TITLE ====================
\begin{center}
{\LARGE\bfseries\color{sectionblue} CSE 711 — Compiler}\\[4pt]
{\large\bfseries All Non-Numerical (Theory) Questions — Model Answers}\\[3pt]
{\small Past papers 2016--2024 \quad$\cdot$\quad University of Chittagong \quad$\cdot$\quad Exam: Wed 24 Jun 2026}
\end{center}

\begin{notesbox}
\small
\textbf{Scope.} This sheet answers every \emph{conceptual / written-answer} question (definitions, comparisons,
explanations) found across the 8-year question map. It deliberately \emph{excludes} the construction \&
computation problems (derivations, ambiguity demos, left-recursion elimination, left-factoring, FIRST/FOLLOW,
LL(1)/LR table construction, parse traces, annotated parse trees, dependency graphs, DAGs, TAC/quad/triple
generation, activation-record diagrams, CFG construction, instruction-cost calculation) — those need worked
diagrams, not prose. Year tags \yr{20XX} show where each concept was asked.
\end{notesbox}

\tableofcontents
\vspace{6pt}

%==================== TOPIC 1 ====================
\section{Lexical Analysis}

\subsection{Token, Pattern, Lexeme \yr{2024 Q2c}\,\yr{2022 A-2}\,\yr{2021 Q3a}\,\yr{2018 Q2a}\,\yr{2017 Q1b}\,\yr{2016 Q2a}}
\begin{qbox}
\begin{itemize}
\item \term{Token:} a pair \(\langle\) token-name, optional attribute value \(\rangle\). The token-name is an abstract
symbol (a terminal of the grammar) the parser works with, e.g.\ \texttt{id}, \texttt{num}, \texttt{if}, \texttt{relop}.
\item \term{Pattern:} the rule — usually a regular expression — that describes the \emph{set} of strings a token may stand for.
\item \term{Lexeme:} the actual sequence of characters in the source program that matches a pattern and is recognised as one instance of a token.
\end{itemize}
\textbf{Relationship:} the scanner reads characters and forms a \emph{lexeme}; the lexeme matches a \emph{pattern};
the pattern defines a \emph{token} that is passed to the parser.

\medskip
\begin{tabularx}{\linewidth}{@{}l l Y@{}}
\toprule
\textbf{Token} & \textbf{Pattern} & \textbf{Sample lexemes}\\
\midrule
\texttt{id} & letter followed by letters/digits & \texttt{count}, \texttt{x1}, \texttt{position}\\
\texttt{num} & digit$^+$ (with optional fraction/exp) & \texttt{27}, \texttt{3.14}\\
\texttt{relop} & \texttt{< | <= | = | <> | > | >=} & \texttt{<=}, \texttt{>}\\
\texttt{if} & the characters \texttt{i f} & \texttt{if}\\
\bottomrule
\end{tabularx}
\end{qbox}

\subsection{Scanning vs. Parsing; Lexer/Parser Tools \yr{2024 Q2d}\,\yr{2023 A-4a}}
\begin{qbox}
\begin{tabularx}{\linewidth}{@{}l Y Y@{}}
\toprule
 & \textbf{Scanning (Lexical Analysis)} & \textbf{Parsing (Syntax Analysis)}\\
\midrule
Job & groups characters into tokens; strips whitespace/comments & groups tokens into grammatical phrases; builds the parse tree\\
Governed by & regular expressions / DFA & context-free grammar / pushdown automaton\\
Power & regular languages & context-free languages\\
Output & stream of tokens & parse / syntax tree\\
Errors caught & illegal characters, malformed tokens & wrong token order, missing \texttt{)}, \texttt{;}\\
\bottomrule
\end{tabularx}

\smallskip
\textbf{Tools.} \term{LEX / Flex} generates a scanner from regular-expression rules; \term{YACC / Bison}
generates a parser from a context-free grammar with semantic actions. They are used together: LEX feeds tokens to the YACC-generated parser.
\end{qbox}

\subsection{Symbol Table — Role, Components, Hash Implementation \yr{2023 A-1c}}
\begin{qbox}
\textbf{Role.} A data structure that records information about every identifier (name) in the program. It is built by
the lexer/parser and consulted by \emph{every} later phase for scope resolution, type checking, and address/code generation.

\textbf{Components of an entry:}
\begin{itemize}
\item name (the lexeme), token/kind (variable, function, type\dots)
\item data type and storage class / scope
\item size and memory offset / address
\item line of declaration; for functions: parameter list \& return type
\end{itemize}
\textbf{Hash-table implementation (C).} An array of buckets indexed by a hash of the name; collisions resolved by
\emph{chaining} (each bucket is a linked list of entries). \texttt{insert} and \texttt{lookup} are $O(1)$ on average.
\end{qbox}

\subsection{Input Buffering \yr{2017 Q1b}}
\begin{qbox}
To read source efficiently the scanner uses a \term{two-buffer scheme}: two halves of size $N$ filled by one system read each.
Two pointers — \texttt{lexemeBegin} (start of current lexeme) and \texttt{forward} (scans ahead). A \term{sentinel} (a
special \texttt{eof} char) is placed at the end of each half so a \emph{single} test on \texttt{forward} detects both
``end of buffer'' and ``end of file,'' minimising comparisons per character.
\end{qbox}

%==================== TOPIC 2 ====================
\section{Top-Down Parsing}

\subsection{Top-Down vs. Bottom-Up Parsing \yr{2024 Q4a}}
\begin{qbox}
\begin{tabularx}{\linewidth}{@{}l Y Y@{}}
\toprule
 & \textbf{Top-Down} & \textbf{Bottom-Up}\\
\midrule
Tree built & root $\to$ leaves & leaves $\to$ root\\
Derivation & \emph{leftmost} & \emph{rightmost}, in reverse\\
Technique & predictive / recursive-descent (LL) & shift-reduce (LR family)\\
Decision & expands a nonterminal using lookahead & reduces a \emph{handle} on the stack\\
Left recursion & not allowed (infinite loop) & no problem\\
Grammar class & smaller (LL $\subset$ LR) & larger / more powerful\\
\bottomrule
\end{tabularx}
\end{qbox}

\subsection{Recursive-Descent Parsing — How It Works \yr{2024 Q4b}\,\yr{2022 A-4}\,\yr{2016 Q3a}}
\begin{qbox}
\begin{itemize}
\item One \emph{procedure per nonterminal}. The procedure body mirrors the right-hand side: it \emph{matches} terminals
against the current input and \emph{calls} the procedure of each nonterminal it meets.
\item \term{General} recursive descent may need \emph{backtracking} (try an alternative, undo input on failure).
\item \term{Predictive} recursive descent uses one lookahead token to choose the alternative with \emph{no} backtracking
— it requires a grammar that is left-factored and free of left recursion.
\item \textbf{Issue:} a left-recursive production $A\to A\alpha$ makes the procedure call itself forever — left recursion must be eliminated first.
\end{itemize}
A procedure typically uses helpers \texttt{match()/scan()} (consume expected token) and \texttt{error()} (report \& recover).
\end{qbox}

\subsection{Ambiguity \yr{2024 Q3d}\,\yr{2022 A-3}}
\begin{qbox}
A grammar is \term{ambiguous} if some string in its language has \emph{more than one} parse tree (equivalently more than
one leftmost or more than one rightmost derivation). It is a problem because the parser cannot decide which structure —
hence which meaning — to assign. Classic example: \(E\to E+E \mid E*E \mid (E) \mid \texttt{id}\) gives two trees for
\texttt{id+id*id}. Fixes: rewrite with precedence/associativity levels, or impose a disambiguating rule.
\end{qbox}

\subsection{Dangling-Else Problem \yr{2023 A-4b}}
\begin{qbox}
The grammar
\[ \texttt{stmt}\to \texttt{if } E \texttt{ then stmt} \mid \texttt{if } E \texttt{ then stmt else stmt}\mid \texttt{other} \]
is ambiguous: in \texttt{if }$E_1$\texttt{ then if }$E_2$\texttt{ then }$S_1$\texttt{ else }$S_2$ the \texttt{else}
can bind to either \texttt{if}. \textbf{Rule used in practice:} match each \texttt{else} with the \emph{nearest unmatched}
\texttt{if}. The disambiguated grammar separates \emph{matched} and \emph{open} statements:
\[
\begin{aligned}
\texttt{stmt} &\to \texttt{matched} \mid \texttt{open}\\
\texttt{matched} &\to \texttt{if } E \texttt{ then matched else matched} \mid \texttt{other}\\
\texttt{open} &\to \texttt{if } E \texttt{ then stmt} \mid \texttt{if } E \texttt{ then matched else open}
\end{aligned}
\]
which forces every \texttt{else} onto the closest preceding \texttt{then}, removing the ambiguity.
\end{qbox}

\subsection{Error-Recovery Strategies \yr{2023 A-3b}\,\yr{2022 A-2}\,\yr{2020 Q1c}\,\yr{2018 Q1a}}
\begin{qbox}
\begin{itemize}
\item \term{Panic mode:} on an error, skip input tokens until a \emph{synchronising} token (e.g.\ \texttt{;}, \texttt{)}, \texttt{end}) appears, then resume. Simple, never loops.
\item \term{Phrase-level:} perform a local correction on the remaining input (replace/insert/delete a few symbols) to continue.
\item \term{Error productions:} add grammar rules for common mistakes so the parser can recognise and diagnose them.
\item \term{Global correction:} find the \emph{minimum} number of changes to turn the input into a valid string — theoretically ideal but too costly in practice.
\end{itemize}
\textbf{In predictive (LL) parsing:} use FOLLOW($A$) tokens to \emph{synchronise} (pop $A$ and resume) and FIRST($A$) tokens to decide what to skip.
\end{qbox}

\subsection{LL vs. LR Parser \yr{2022 B-1a}}
\begin{qbox}
\begin{tabularx}{\linewidth}{@{}l Y Y@{}}
\toprule
 & \textbf{LL} & \textbf{LR}\\
\midrule
Scan & \underline{L}eft-to-right & \underline{L}eft-to-right\\
Derivation traced & \underline{L}eftmost & \underline{R}ightmost (reverse)\\
Direction & top-down (predictive) & bottom-up (shift-reduce)\\
Left recursion & cannot handle & handles it\\
Power & weaker: LL(1) $\subset$ LR(1) & more grammars accepted\\
Table & predictive parsing table & ACTION + GOTO tables\\
\bottomrule
\end{tabularx}
\end{qbox}

\subsection{Why Left Recursion Is Not a Problem for Bottom-Up Parsing \yr{2022 A-3}}
\begin{qbox}
A shift-reduce parser builds the tree from the \emph{leaves up}: it shifts terminals onto the stack and reduces a
\emph{handle} the moment one appears on top. For $A\to A\alpha$ it simply shifts the leading symbols and reduces left-to-right,
so the recursion terminates naturally. A top-down parser, by contrast, would re-enter $A$ before reading any input and
loop forever. Hence left recursion is harmless (even idiomatic) for bottom-up parsers, but fatal for top-down ones.
\end{qbox}

\subsection{Suitability for Top-Down / LL(1) \yr{2023 A-2a}}
\begin{qbox}
A grammar is fit for predictive top-down (LL(1)) parsing iff for every nonterminal $A$ with alternatives
$A\to\alpha\mid\beta$:
\begin{enumerate}
\item it has \emph{no left recursion} and is \emph{left-factored};
\item $\text{FIRST}(\alpha)\cap\text{FIRST}(\beta)=\varnothing$; and
\item if $\beta\Rightarrow^*\varepsilon$ then $\text{FIRST}(\alpha)\cap\text{FOLLOW}(A)=\varnothing$.
\end{enumerate}
These conditions guarantee a unique production choice from one lookahead token.
\end{qbox}

%==================== TOPIC 3 ====================
\section{Bottom-Up Parsing}

\subsection{LR(0) vs. SLR(1) vs. LR(1) vs. LALR(1) \yr{2024 Q6a}}
\begin{qbox}
\begin{tabularx}{\linewidth}{@{}l Y@{}}
\toprule
\textbf{Parser} & \textbf{Characteristics}\\
\midrule
\term{LR(0)} & no lookahead; reduces in any state holding $A\to\alpha\bullet$; weakest, most conflicts.\\
\term{SLR(1)} & LR(0) items + reduce by $A\to\alpha$ only when next token $\in$ FOLLOW($A$); small tables, fewer conflicts.\\
\term{LR(1)} (canonical) & items carry an explicit lookahead symbol; most powerful; largest number of states/tables.\\
\term{LALR(1)} & merges LR(1) states with the same core; same \#states as SLR, almost as strong as LR(1); used by YACC. May add reduce–reduce conflicts but never shift–reduce.\\
\bottomrule
\end{tabularx}
\smallskip
\textbf{Power:} \(\text{LR(0)} \subset \text{SLR(1)} \subset \text{LALR(1)} \subset \text{LR(1)}\).
\end{qbox}

\subsection{Shift-Reduce Parsing Conflicts \yr{2023 B-1a}}
\begin{qbox}
\begin{itemize}
\item \term{Shift-reduce conflict:} in some state the parser can neither decide to \emph{shift} the next token nor to
\emph{reduce} a completed production (classic case: the dangling-else).
\item \term{Reduce-reduce conflict:} two different productions are both candidates for reduction in the same state on the same lookahead.
\end{itemize}
Both arise from ambiguous or non-LR grammars; resolving them needs more lookahead, grammar rewriting, or a stated precedence rule.
\end{qbox}

\subsection{Which Parser Class Can Parse a Grammar \yr{2017 Q5a}}
\begin{qbox}
To classify a grammar, test it bottom of the hierarchy upward and report the smallest class that succeeds:
\[ \text{LL(1)} \subset \text{SLR(1)} \subset \text{LALR(1)} \subset \text{LR(1)}. \]
Check LL(1) by the disjoint-FIRST/FOLLOW conditions (needs no left recursion, left-factored). Check the LR family by
building item sets and looking for shift-reduce / reduce-reduce conflicts. An ambiguous grammar belongs to \emph{none}
of these deterministic classes.
\end{qbox}

%==================== TOPIC 4 ====================
\section{Syntax-Directed Translation}

\subsection{Synthesized vs. Inherited Attributes \yr{2024 Q6c}\,\yr{2023 B-1b}\,\yr{2022 B-1b}\,\yr{2021 Q5}\,\yr{2020 Q3b}\,\yr{2018 Q2c}\,\yr{2016 Q2c}}
\begin{qbox}
\begin{tabularx}{\linewidth}{@{}l Y Y@{}}
\toprule
 & \textbf{Synthesized} & \textbf{Inherited}\\
\midrule
Computed from & attributes of the node's \emph{children} (and itself) & attributes of \emph{parent} and/or \emph{siblings}\\
Information flow & bottom-up & top-down / left-to-right\\
Defined at & $A\to\alpha$, value of $A$ from RHS symbols & a symbol on the RHS, value from parent/left siblings\\
Typical use & value of an expression, code of a subtree & passing a declared type down to a list of ids\\
\bottomrule
\end{tabularx}
\smallskip
A \emph{terminal} has only synthesized attributes (supplied by the lexer); the \emph{start symbol} has no inherited attributes.
\end{qbox}

\subsection{S-Attributed vs. L-Attributed Definitions \yr{2024 Q6c}\,\yr{2023 B-3a}\,\yr{2018 Q2c}\,\yr{2017 Q4c}\,\yr{2016 Q2c}}
\begin{qbox}
\begin{itemize}
\item \term{S-attributed SDD:} uses \emph{only synthesized} attributes. It can be evaluated bottom-up during an LR parse
(attributes computed on reduction).
\item \term{L-attributed SDD:} every inherited attribute of a RHS symbol $X_j$ depends only on the inherited attributes
of the parent and on attributes of symbols \emph{to the left} of $X_j$ (plus any synthesized attributes). It can be
evaluated in a single left-to-right, depth-first pass — compatible with LL (predictive) parsing.
\end{itemize}
\textbf{Relationship:} every S-attributed definition is also L-attributed (no inherited attributes to violate the rule), so $\text{S-attributed}\subset\text{L-attributed}$.
\end{qbox}

%==================== TOPIC 5 ====================
\section{Intermediate Code}

\subsection{Quadruples vs. Triples vs. Indirect Triples vs. SSA \yr{2020 Q7c}}
\begin{qbox}
\begin{tabularx}{\linewidth}{@{}l Y@{}}
\toprule
\textbf{Form} & \textbf{Description}\\
\midrule
\term{Quadruple} & \((op,\ arg_1,\ arg_2,\ result)\) — results held in \emph{explicit temporary names}. Easy to move/reorder for optimization; uses most space.\\
\term{Triple} & \((op,\ arg_1,\ arg_2)\) — a result is referenced by the \emph{position (index)} of its triple. Compact, but reordering is hard because positions are the references.\\
\term{Indirect triple} & a separate \emph{list of pointers} to the triples. Optimization reorders only the pointer list; the triples themselves stay fixed — best of both.\\
\term{SSA} & Static Single Assignment: every variable is assigned \emph{exactly once}; \(\phi\)-functions are inserted at control-flow merge points. Greatly simplifies data-flow optimizations.\\
\bottomrule
\end{tabularx}
\end{qbox}

%==================== TOPIC 6 ====================
\section{Runtime Environments}

\subsection{Activation Record \& Runtime Stack \yr{2018 Q7c}\,\yr{2018 Q8b}}
\begin{qbox}
An \term{activation record} (stack frame) is the storage block for \emph{one activation} of a procedure. Typical fields
(top to bottom):
\begin{itemize}
\item temporaries and local data
\item saved machine state (return address, saved registers)
\item \term{access (static) link} — to the enclosing scope, for non-local names
\item \term{control (dynamic) link} — to the caller's activation record
\item actual parameters and the returned value
\end{itemize}
The \term{runtime stack} holds these records: one is \emph{pushed} on each call and \emph{popped} on return. Because every
call gets a fresh record, the stack naturally supports \emph{recursion}.
\end{qbox}

\subsection{Static vs. Dynamic Storage \& Garbage Collection \yr{2023 B-4}}
\begin{qbox}
\begin{itemize}
\item \term{Static allocation:} storage sizes/addresses fixed at compile time (globals, static data). No recursion, no run-time sizing.
\item \term{Stack (dynamic) allocation:} activation records pushed/popped in LIFO order; supports recursion and block scope.
\item \term{Heap allocation:} for data whose lifetime is unrelated to the call that created it (objects, \texttt{malloc}/\texttt{new}); allocated/freed in any order.
\end{itemize}
\term{Garbage collector:} reclaims heap storage that is no longer reachable from the program, automatically. Common
techniques: \emph{reference counting}, \emph{mark-and-sweep}, and \emph{copying} collection.
\end{qbox}

\subsection{Environment, State, and Static Binding \yr{2017 Q1a}}
\begin{qbox}
\begin{itemize}
\item \term{Environment:} a mapping from a \emph{name} to a \emph{storage location} (its l-value).
\item \term{State:} a mapping from a \emph{storage location} to the \emph{value} held there (its r-value).
\end{itemize}
An assignment changes the \emph{state}; entering a new scope or declaring a name changes the \emph{environment}.
\term{Static (compile-time) binding} fixes an association — such as scope or type — when the program is compiled;
\emph{dynamic binding} fixes it at run time.
\end{qbox}

%==================== TOPIC 7 ====================
\section{Code Optimisation}

\subsection{Basic Block \& Flow Graph \yr{2024 Q7c}}
\begin{qbox}
\begin{itemize}
\item \term{Basic block:} a maximal sequence of consecutive instructions with a \emph{single entry} (the first instruction)
and a \emph{single exit} (the last) — control enters only at the top and leaves only at the bottom, no internal branching.
\item \term{Leaders} (first instruction of each block): the first instruction of the program; any jump target; and any
instruction immediately following a jump.
\item \term{Flow graph (CFG):} a directed graph whose nodes are basic blocks and whose edges show possible flow of control between them.
\end{itemize}
\end{qbox}

\subsection{Objectives of Optimization \yr{2024 B-3b}}
\begin{qbox}
Transform the intermediate/target code so it \emph{runs faster and/or uses less memory}, while:
\begin{itemize}
\item \term{preserving the program's meaning} (semantics-preserving — the output must be identical);
\item being worthwhile \emph{on average} (the speed-up justifies compile-time cost);
\item not degrading the common case to help a rare one.
\end{itemize}
\end{qbox}

\subsection{Copy Propagation \yr{2023 B-2c}}
\begin{qbox}
After a copy statement \(x = y\), replace later \emph{uses} of \(x\) by \(y\) wherever \(x\) is not redefined in between.
\textbf{Advantage:} it often makes the copy statement \emph{dead} (so dead-code elimination can delete it) and exposes
further optimizations such as constant folding and common-subexpression elimination.
\end{qbox}

\subsection{Optimization Techniques (names \& meaning) \yr{2024 B-3b}\,\yr{2016 Q8}}
\begin{qbox}
\begin{tabularx}{\linewidth}{@{}l Y@{}}
\toprule
\textbf{Technique} & \textbf{What it does}\\
\midrule
Common-subexpression elimination & reuse a value already computed instead of recomputing it.\\
Copy propagation & replace uses of $x$ by $y$ after $x=y$.\\
Constant folding / propagation & evaluate constant expressions at compile time; spread known constants.\\
Dead-code elimination & remove code whose result is never used.\\
Strength reduction & replace a costly op by a cheaper one (e.g.\ \(x*2\to x{+}x\), multiply $\to$ add in loops).\\
Code motion (loop-invariant hoisting) & move computations that don't change inside a loop to before the loop.\\
Induction-variable elimination & remove redundant loop counters derived from one another.\\
Algebraic simplification & apply identities such as $x+0=x$, $x*1=x$.\\
\bottomrule
\end{tabularx}
\smallskip
These are applied \emph{locally} (within a basic block, often via a DAG) and \emph{globally} (across the flow graph, using data-flow analysis).
\end{qbox}

%==================== TOPIC 8 ====================
\section{Code Generation}

\subsection{Stack Frames; Caller- vs. Callee-Saved Registers \yr{2021 Q7c}}
\begin{qbox}
\textbf{Stack frame.} The per-call region on the runtime stack holding the return address, saved registers, local
variables and spilled temporaries; addressed via a \emph{frame pointer} (\texttt{rbp}) with the \emph{stack pointer}
(\texttt{rsp}) at the top.

\smallskip
\begin{tabularx}{\linewidth}{@{}l Y@{}}
\toprule
\textbf{Register class} & \textbf{Who preserves it}\\
\midrule
\term{Caller-saved} (volatile) & the \emph{caller} must save them before a call if it needs the values afterward; the callee may clobber them freely. \emph{(x86-64 SysV: \texttt{rax}, \texttt{rcx}, \texttt{rdx}, \texttt{rsi}, \texttt{rdi}, \texttt{r8–r11}.)}\\
\term{Callee-saved} (non-volatile) & the \emph{callee} must save and restore them, so the caller may assume they survive the call. \emph{(\texttt{rbx}, \texttt{rbp}, \texttt{r12–r15}.)}\\
\bottomrule
\end{tabularx}
\end{qbox}

\vfill
\begin{center}\small\color{sectionblue!70!black}
End of theory model answers — CSE 711 Compiler (2016–2024). Good luck.
\end{center}

\end{document}
